% Generated draft: requires mathematical and source review.
\documentclass[11pt]{article}
\usepackage[margin=25mm]{geometry}
\usepackage{fontspec}
\setmainfont{Latin Modern Roman}
\usepackage{amsmath,amssymb,mathtools,unicode-math}
\setmathfont{Latin Modern Math}
\usepackage{enumitem,needspace,xurl,hyperref}
\setlength{\parindent}{0pt}
\setlength{\parskip}{5pt}
\setlength{\emergencystretch}{3em}
\newcommand{\N}{\mathbb{N}}
\newcommand{\Z}{\mathbb{Z}}
\newcommand{\Q}{\mathbb{Q}}
\newcommand{\R}{\mathbb{R}}
\newcommand{\C}{\mathbb{C}}
\newcommand{\F}{\mathbb{F}}
\newcommand{\A}{\mathbb{A}}
\newcommand{\PP}{\mathbb{P}}
\newcommand{\E}{\mathbb{E}}
\newcommand{\eps}{\varepsilon}
\newcommand{\dd}{\,\mathrm{d}}
\begin{document}
\section*{UnsolvedMath ID: 30002973}
\textbf{Source problem code:} catalogue entry
\subsection{Definitions and mathematical statement}
A complex fourfold \(X\) is of \(K3^{[2]}\)-type if it is an irreducible holomorphic symplectic manifold deformation equivalent to the Hilbert scheme of length-two subschemes of a complex \(K3\) surface. Here irreducible holomorphic symplectic means compact, simply connected, and K\"ahler, with \(H^0(X,\Omega_X^2)=\C\sigma\), where \(\sigma\) is everywhere nondegenerate. An automorphism \(g\in\operatorname{Aut}(X)\) is nonsymplectic of prime order \(p\) if \(\operatorname{ord}(g)=p\) and \(g^*\sigma=\zeta\sigma\) for a primitive \(p\)-th root of unity \(\zeta\ne1\).

For \(d\geq1\), the complex ball is
\[
\mathbb B^d=\left\{(z_1,\ldots,z_d)\in\C^d:\sum_{j=1}^d|z_j|^2<1\right\}.
\]
A complex ball quotient has the form \(\Gamma\backslash\mathbb B^d\), with \(\Gamma\) a discrete group of holomorphic automorphisms acting properly discontinuously. The supplied record concerns moduli spaces associated with fourfolds \(X\) equipped with such an automorphism \(g\), in cases admitting the relevant ball quotient description. It does not specify the primes, deformation types of pairs, lattice data, quotient groups, or connected components that define the intended family. Nor does it specify whether the moduli space is the full quotient or an open subset obtained by removing loci from it.

For an irreducible complex quasi-projective variety \(M\), the ordinary birational interpretation of its Kodaira dimension is \(\kappa(Y)\), where \(Y\) is a smooth projective model of its function field. For such a model, \(\kappa(Y)=-\infty\) if \(H^0(Y,mK_Y)=0\) for every integer \(m\geq1\); otherwise it is the maximum dimension of the images of the pluricanonical rational maps associated with the nonzero systems \(|mK_Y|\). This is independent of the smooth projective model. The variety \(M\) is rational over \(\C\) if it is birational over \(\C\) to \(\PP^{\dim M}_{\C}\), equivalently if
\[
\C(M)\cong\C(t_1,\ldots,t_{\dim M})
\]
as extensions of \(\C\).

\textbf{Requested classification, subject to unresolved scope.} For every irreducible moduli space \(M\) in the source's intended family of complex ball quotient moduli spaces associated with \(K3^{[2]}\)-type fourfolds carrying a nonsymplectic automorphism of prime order, determine its Kodaira dimension and decide whether it is rational over \(\C\). Under the ordinary birational interpretation just defined, the first target is
\[
\kappa(Y)\in\{-\infty,0,1,\ldots,\dim M\}.
\]
The supplied record does not determine the family precisely enough to enumerate the spaces quantified over. It also does not distinguish ordinary Kodaira dimension from logarithmic Kodaira dimension for a noncompact moduli space. Thus the definition above makes one conventional interpretation explicit without asserting that the source selects it. The original wording does not restrict the automorphism order to a prime, whereas the supplied clean statement does; this draft retains the clean statement's prime-order restriction and records the discrepancy for review.
\subsection{Short English statement}
Determine the Kodaira dimensions and rationality of the intended complex ball quotient moduli spaces of \(K3^{[2]}\)-type fourfolds with a nonsymplectic automorphism of prime order. The exact family and the intended interpretation of Kodaira dimension require clarification. The supplied status is partially solved, based on dated curation rather than independent verification.
\subsection{Sources}
\begin{itemize}
\item Local UnsolvedMath registry: permanent ID \(30002973\), source problem code \texttt{catalogue entry}, title \emph{Birational Geometry of Ball-Quotient Moduli Spaces}, in the Oberwolfach Reports Open Problems collection. No individual proposer is supplied. The proposed year is recorded as \(2015\).
\item Supplied report citation: \emph{Mini-Workshop: Singular Curves on \(K3\) Surfaces and Hyperk\"ahler Manifolds} (2016), \emph{Oberwolfach Reports}, pp.~2939--2967, DOI \texttt{10.4171/owr/2015/51}. \url{https://doi.org/10.4171/owr/2015/51}. The citation's year \(2016\), the recorded proposed year \(2015\), and the DOI's \(2015\) component are retained as supplied; their relationship has not been independently checked.
\item Supplied literature reference: Boissiere--Camere--Sarti, \emph{Complex Ball Quotients}. \url{https://doi.org/10.1016/j.jalgebra.2018.03.012}.
\item Supplied literature reference: Ma--Ohashi--Taki, \emph{Rationality of Eisenstein K3 Moduli}. \url{https://arxiv.org/abs/1208.5826}.
\item Supplied literature reference: \emph{Oberwolfach Ball-Quotient Questions}. \url{https://oa.tib.eu/renate/bitstreams/205e22d5-f548-4915-b721-7e608401c35b/download}.
\item Literature assessment dated \texttt{2026-08-21}: the supplied curation labels the problem \emph{partially solved}, states that the period-map construction of the relevant ball quotients is established and many order-three \(K3\) cases are rational, and states that Kodaira dimensions and rationality of the full higher-dimensional family remain unsettled. This is supplied dated curation, not an independently verified literature assessment. Its claim about \(K3\) surfaces does not itself establish a corresponding claim about fourfold moduli.
\item The registry labels the statement \texttt{corrected\_verified} and reports that it was checked against the official Oberwolfach report and preserves its hypotheses, notation, and mathematical target. These are supplied verification claims; no source visit or independent verification is asserted here.
\end{itemize}
\end{document}
